\documentclass{openjournal}

\usepackage[T1]{fontenc}
\usepackage{graphicx}
\usepackage{amsmath,amssymb}
\usepackage{hyperref}

\shorttitle{A Volume-Limited ALMA Search for Technosignatures within 20\,pc}
\shortauthors{White \& Dey}

\begin{document}

\title{A Volume-Limited Search for Technosignatures within 20\,pc Using the Public ALMA Archive: I. Background, Motivation and Survey Design}

\author{Glenn J. White\altaffilmark{1,2} and Robin Dey\altaffilmark{3}}
\altaffiltext{1}{School of Physical Sciences, The Open University, Walton Hall, Milton Keynes, MK7 6AA, UK; glenn.white@open.ac.uk}
\altaffiltext{2}{RAL Space, STFC Rutherford Appleton Laboratory, Harwell Campus, Didcot, OX11 0QX, UK}
\altaffiltext{3}{VBRL Holdings Inc.}

\begin{abstract}
We describe the motivation, background and design of a volume-limited search for
radio technosignatures among all stars within 20\,pc of the Sun that have public
observations in the Atacama Large Millimeter/submillimeter Array (ALMA) science
archive. Unlike the great majority of technosignature surveys conducted to date,
which are either wide, shallow all-sky surveys or targeted campaigns preferentially
selecting exoplanet hosts, young stellar objects, or other astrophysically
``interesting'' systems, our sample is constructed to be representative of the
general stellar population in the immediate solar neighbourhood, irrespective of
spectral type, age, multiplicity, or known planetary or circumstellar-disc status.
We mine the ALMA archive opportunistically: rather than requesting dedicated
observing time, we reprocess existing calibrated (and, where necessary,
recalibrated) visibility data to search simultaneously for narrowband
Doppler-drifting technosignature candidates and for anomalous broad-band continuum
emission, at frequencies ($\sim$\,90--950\,GHz) that remain almost entirely
unexplored by the technosignature literature. This paper (Paper~I of a series)
reviews six decades of radio technosignature searches, the small number of
millimetre/submillimetre efforts to date, recent applications of machine learning
to technosignature detection, and the statistical frameworks used to interpret
non-detections, in order to place the present survey in context and to motivate
its volume-limited, archival design. The observational results of the survey
itself, together with the resulting upper limits on transmitter prevalence in the
solar neighbourhood, will be presented in Paper~II.
\end{abstract}

\keywords{extraterrestrial intelligence --- radio lines: stars --- submillimetre: stars --- surveys --- astrobiology}

\section{Introduction}
\label{sec:intro}

The search for extraterrestrial intelligence (SETI) rests on the premise that a
technological civilisation, whether deliberately signalling or simply going about
its own affairs, may leave a detectable imprint on the electromagnetic spectrum: a
\emph{technosignature} \citep{Sheikh2020,SETIreview2025}. Since the foundational
proposal of \citet{CocconiMorrison1959} that microwave frequencies near the 21\,cm
hydrogen line would be an obvious ``meeting place'' for interstellar communication,
and Frank Drake's first observational test of the idea with Project Ozma
\citep{Drake1960}, the radio band has remained the dominant search space for
technosignatures, for reasons that are both historical and physically well
motivated: radio photons are cheap to produce in enormous numbers, propagate
through the interstellar medium and planetary atmospheres with little attenuation,
and (for a sufficiently narrowband signal) can in principle be distinguished from
any known natural astrophysical process by their spectral compactness and Doppler
dynamics \citep{Sheikh2020}.

Six decades on, radio technosignature searches have grown from single-target,
single-frequency pointings into surveys of many thousands of stars across broad
swathes of the radio spectrum. Yet, as \citet{Wright2018} demonstrated
quantitatively with their ``Cosmic Haystack'' formalism, even the combined efforts
of every major SETI programme to date have sampled only an infinitesimal fraction
of the plausible multi-dimensional parameter space of transmitter frequency,
bandwidth, power, location, and duty cycle. Persistent, non-trivial gaps remain:
almost the entire millimetre and submillimetre regime above $\sim$\,50\,GHz is
essentially unsearched; most surveys have targeted specific classes of star
(nearby exoplanet hosts, young stellar objects with protoplanetary discs, stars in
the ``Earth Transit Zone'') rather than an unbiased, complete sample of the local
stellar population; and only a small number of studies have systematically
reprocessed \emph{existing} interferometric archives for technosignatures, as
opposed to requesting dedicated telescope time.

This paper is the first in a series describing a new survey designed to address
several of these gaps simultaneously. We construct a volume-limited sample of
every star within 20\,pc of the Sun with public ALMA archival coverage,
irrespective of spectral type, disc-hosting status, or known planets, and mine
that archival data for both narrowband technosignature candidates and anomalous
continuum emission, applying corrections for stellar proper motion and for the
range of Doppler drift rates plausible for a transmitter on a planet in orbit
around its host star. Section~\ref{sec:background} reviews the relevant
literature: the history of radio SETI (\S\ref{sec:history}), the current
generation of large-scale technosignature surveys (\S\ref{sec:modern}), the very
limited body of work at millimetre/submillimetre frequencies
(\S\ref{sec:mmsubmm}), the growing use of machine learning and anomaly detection
in this field (\S\ref{sec:ml}), the statistical frameworks used to interpret
survey non-detections (\S\ref{sec:stats}), the case for a volume-limited,
representative sample (\S\ref{sec:volumelimited}), and the prospects offered by
next-generation facilities (\S\ref{sec:future}). Section~\ref{sec:sample}
summarises our sample selection, Section~\ref{sec:method} our data-reduction and
search methodology, and Section~\ref{sec:futurework} the scope of the results to
be presented in Paper~II.

\section{Background}
\label{sec:background}

\subsection{Six decades of radio SETI: a brief history}
\label{sec:history}

Project Ozma \citep{Drake1960} observed two nearby Sun-like stars, $\tau$~Ceti and
$\epsilon$~Eridani, with the 26\,m Tatel telescope at Green Bank at 1420\,MHz for
a total of some 150 hours, and found no signal, but established both the
methodology and the symbolic starting point for the modern field. Through the
1960s--1980s a steady stream of targeted and sky-survey searches followed,
including Ozma~II \citep{Drake1986}, the Ohio State ``Big Ear'' survey (source of
the famous ``Wow!'' signal), and the META and BETA all-sky surveys conducted from
Harvard/Smithsonian facilities. Project Phoenix \citep{Backus2002}, run by the
SETI Institute from 1995--2004, was the most sensitive targeted survey of its era,
observing on the order of 800 nearby Sun-like stars using the Arecibo, Parkes and
Green Bank telescopes, and pioneered a rigorous methodology for distinguishing
candidate signals from terrestrial radio-frequency interference (RFI) through
simultaneous multi-site or multi-antenna verification. The Allen Telescope Array
\citep[ATA;][]{Tarter2001}, jointly developed by the SETI Institute and UC
Berkeley and first commissioned in 2007, was the first radio telescope
purpose-built from the outset for commensal, wide-band SETI observing, and remains
in active use, including in recent technosignature searches of interstellar
objects \citep{ATAISO2025} and TRAPPIST-1 \citep{ATATRAPPIST2024}. The
distributed-computing project SETI@home \citep{SETIhome2025} processed
public-participant-donated CPU cycles to search Arecibo sky-survey data for
narrowband and pulsed signals for two decades before its 2020 shutdown, and its
final data-analysis summary was only published in 2025.

Two complementary observing strategies have been used throughout this history:
\emph{sky surveys}, covering a large fraction of the celestial sphere at modest
sensitivity per pointing, and \emph{targeted searches}, concentrating integration
time on a preselected list of promising stars \citep[e.g.][]{Drake1974}. Both
strategies persist in the modern era, but -- as we discuss in
\S\ref{sec:volumelimited} -- the majority of targeted programmes have selected
their samples on some presumed indicator of habitability or interest (proximity,
Sun-like spectral type, known planets, or membership of the Earth Transit Zone),
rather than aiming for a complete, unbiased census of the local stellar
population.

\subsection{Modern large-scale radio technosignature surveys}
\label{sec:modern}

The current generation of technosignature science is dominated by the
Breakthrough Listen initiative, launched in 2015 with a ten-year, \$100\,million
commitment to survey the one million nearest stars, the full Galactic plane and
centre, the 100 nearest galaxies, and a sample of exoplanet and planet-candidate
hosts, primarily using the Robert C.\ Byrd Green Bank Telescope (GBT) and the
CSIRO Parkes ``Murriyang'' telescope, and now increasingly via commensal
observing on MeerKAT \citep{BLoverview}. Its first major published dataset,
\citet{Enriquez2017}, searched 692 nearby stars over 1.1--1.9\,GHz with the GBT
and set limits on narrowband transmitters of EIRP\,$>10^{13}$\,W; this was
followed by \citet{Price2020} (1327 stars, 1.10--3.45\,GHz),
\citet{Gajjar2019EarthTransitZone}-style searches of the restricted Earth Transit
Zone at 3.95--8.00\,GHz, and, most recently, a substantial expansion of automated,
commensal technosignature searching with MeerKAT that has processed some 1.5
million coherent beams between 2023 and 2026 \citep{BLMeerKAT2026}. In parallel,
\citet{Margot2023} conducted a large citizen-science-style search of 11\,680
stars with the GBT at 1.15--1.73\,GHz, explicitly developing a statistical
formalism for interpreting non-detections in terms of the fraction of stars that
could host a detectable transmitter, \emph{provided the observed sample is
representative of the underlying stellar population of the search volume} -- a
condition that, as we discuss below, is often not satisfied by targeted,
interest-selected samples.

Beyond Breakthrough Listen, the last five years have seen international radio
facilities extend technosignature searches to new frequency regimes and
observational strategies: the Five-hundred-metre Aperture Spherical Telescope
\citep[FAST;][]{Zhang2020, FAST3IATLAS2025} in China has conducted commensal
searches with unmatched raw sensitivity at L-band, including recent narrowband
searches of interstellar objects; LOFAR and the pathfinder NenuFAR have pushed
technosignature searches down to 10--190\,MHz, with LOFAR alone having surveyed
over 1.6 million stars in commensal mode and dual-site LOFAR configurations
enabling near-real-time RFI rejection \citep{LOFARdualsite2023}; the Sardinia
Radio Telescope has conducted one of the first dedicated high-frequency
($>$\,18\,GHz) surveys of nearby stars and the Galactic Centre
\citep{SRT2025}; and commensal technosignature back-ends are now operating or in
development at the Karl G.\ Jansky Very Large Array \citep[the COSMIC
system;][]{Tremblay2024}, MeerKAT \citep{Czech2021}, and the Murchison Widefield
Array \citep{Morrison2023}. Occultation-based strategies -- observing a star
during a planet-planet or planet-star occultation, when geometric alignment
transiently favours detection of leakage or deliberate signals -- have also grown
in popularity, exemplified by recent Breakthrough Listen searches of TESS-selected
eclipsing systems \citep{TESSeclipsing2025}.

\subsection{Millimetre and submillimetre technosignature searches}
\label{sec:mmsubmm}

In stark contrast to the radio regime below $\sim$\,10\,GHz, the millimetre and
submillimetre bands remain almost entirely unexplored for technosignatures.
\citet{Mason2024} conducted, to our knowledge, the first technosignature search
ever performed with ALMA, searching archival Band\,3 data (spectral windows
centred at 90.642 and 93.151\,GHz) for narrowband signals towards 28 Galactic
stars selected from Gaia~DR3. Their approach used a ``stellar bycatch'' method:
rather than targeting stars directly, they searched for serendipitous stars
falling within the undistorted field of view of four bright ALMA phase
calibrators, since these calibrators are observed frequently and are not
generally themselves of astrophysical interest. For the nearest star in their
sample they placed a limit of EIRP$_{\rm min} > 7\times10^{17}$\,W, and discussed
in detail the specific challenges of high-frequency technosignature work:
substantially larger Doppler drift rates at fixed line-of-sight acceleration
(since drift rate scales linearly with observing frequency), the resulting
sensitivity loss (``smearing'') for a fixed spectral and temporal resolution
unless a drift-rate search is performed, and an increased incidence of spectral
confusion from the dense molecular line forest present in protoplanetary discs
and star-forming regions at these frequencies. \citet{ATATRAPPIST2024} extended
technosignature work on the TRAPPIST-1 system to the Allen Telescope Array,
underlining the continued observational interest in this particular target but
also its status as one of very few systems that has now been searched for
technosignatures across a wide range of frequencies. Our own earlier work (White
\& Dey, in prep.; archived at
\url{https://github.com/Tilanthi/SETI}) reproduced and extended the
\citet{Mason2024} methodology directly on ALMA science-target (rather than
calibrator-bycatch) fields for TRAPPIST-1 across Bands 3, 6 and 7, finding no
detection and placing limits of EIRP$_{\rm min}\sim5$--$10\times10^{13}$\,W, and
subsequently built a ranked target list and pilot survey of 100 nearby,
ALMA-observed stars, from which the present, volume-limited survey has grown.
Beyond these efforts, we are not aware of any other published, systematic
technosignature search using ALMA or any comparably sensitive
millimetre/submillimetre interferometer, representing a substantial and readily
addressable gap in frequency coverage relative to the sub-10\,GHz literature.

The near-complete absence of mm/submm technosignature work to date is
attributable less to any fundamental unsuitability of these frequencies -- Sheikh's
``nine axes of merit'' framework \citep{Sheikh2020} notes that higher frequencies
suffer less from interstellar scintillation and dispersion, and offer access to
much larger instantaneous bandwidths -- than to historical circumstance:
sensitive mm/submm interferometers with the necessary spectral resolution and
archival depth (chiefly ALMA) have only existed, and only accumulated a
sufficiently large calibrated science archive, within roughly the last decade.
This makes archival reprocessing, rather than dedicated proposal-driven
observing, a particularly natural strategy in this regime: ALMA's archive already
contains many thousands of hours of high-spectral-resolution observations of
individual, nearby stars, taken for entirely unrelated astrophysical purposes
(disc structure, molecular chemistry, stellar activity, debris-disc detection,
calibration monitoring), each of which can in principle be re-searched for
technosignatures at negligible additional observing cost.

\subsection{Machine learning and anomaly detection in technosignature searches}
\label{sec:ml}

The scale of modern technosignature datasets -- hundreds of hours of
high-time- and frequency-resolution data per survey, often containing many
thousands of candidate ``hits'' dominated by anthropogenic RFI -- has made
machine learning an increasingly central tool for both signal detection and RFI
rejection. \citet{Ma2023} presented the first large-scale deep-learning
technosignature search, applying a $\beta$-convolutional variational autoencoder
to 480 hours of Breakthrough Listen GBT data on 820 nearby stars in a
semi-unsupervised anomaly-detection mode, recovering candidate signals missed by
the standard \textsc{turboSETI} narrowband drift-search pipeline
\citep{Enriquez2019} and reducing the number of signals requiring human vetting by
roughly two orders of magnitude, while returning eight signals of interest for
re-observation. Building on this, \citet{Choza2024} identified and corrected a
long-standing sensitivity mischaracterisation in \textsc{turboSETI}'s
signal-to-noise calibration, prompting a retrospective correction of quoted
sensitivities in several earlier Breakthrough Listen publications
\citep{GBTarchiveML2025}. More recently, unsupervised clustering methods (e.g.\
the ``GLOBULAR'' algorithm) have been developed specifically to separate genuine
anomalies from RFI in narrowband technosignature searches without requiring
labelled training examples \citep{GLOBULAR2025}, and further anomaly-detection
searches have been applied to combined Parkes and GBT Breakthrough Listen archives
\citep{AnomalyParkesGBT2025}. These efforts collectively demonstrate that
machine-learning-based anomaly detection can recover classes of signal
(broadband, non-linearly drifting, or otherwise morphologically unusual) that
narrowband, linear-drift matched-filter searches such as \textsc{turboSETI} are,
by design, insensitive to \citep{BRaTs2026}, and that a purely classical,
threshold-based search -- of the kind we adopt for the present survey, following
\citet{Mason2024} -- necessarily leaves some fraction of the true technosignature
parameter space unexamined. We return to this point in \S\ref{sec:futurework} as
a natural direction for follow-up work once the present survey's baseline dataset
exists.

\subsection{Statistical frameworks for interpreting non-detections}
\label{sec:stats}

A non-detection is not, by itself, a null result: it must be translated into a
quantitative statement about which regions of transmitter parameter space have
been excluded, and with what confidence. \citet{Wright2018} formalised this with
the ``Cosmic Haystack'' framework, modelling the technosignature search space as
an eight-dimensional volume (spanning, e.g., transmitter frequency, bandwidth,
location, distance, and duty cycle) and showing that even the combined efforts of
every major SETI programme conducted up to 2018 sampled a small fraction
of any astrophysically plausible haystack volume -- a result frequently
paraphrased as showing that humanity has searched ``the equivalent of a large
hot tub's worth of the Earth's oceans''. \citet{Sheikh2020} subsequently proposed
a complementary ``nine axes of merit'' rubric for comparing the relative
scientific value of different technosignature search strategies (spanning axes
such as detectability, transmitter cost and information content), providing a
qualitative framework alongside Wright et al.'s quantitative one.
\citet{Margot2023} developed an explicit formalism for converting a survey's
non-detection rate into an upper bound on the fraction of stars in the search
volume hosting a detectable transmitter, but noted explicitly that this
inference is only valid if the observed sample is statistically representative
of the underlying stellar population of the search volume -- a requirement that
most targeted, interest-selected samples (including our own earlier 100-star
pilot list) do not, by design, satisfy. Most recently, \citet{EarthDetectingEarth2025}
inverted the usual question, asking at what distance Earth's own technosignatures
would themselves be detectable by an ALMA- or SKA-like observatory, providing a
valuable calibration point (an ``ichnoscale'' of unity) against which the
sensitivity of any given survey, including the present one, can be benchmarked.

\subsection{The case for a volume-limited, representative sample}
\label{sec:volumelimited}

The great majority of the surveys reviewed above -- including our own earlier
100-star ALMA pilot programme -- select targets by some combination of proximity,
Sun-like spectral type, known exoplanets, protoplanetary-disc or debris-disc
detection, or membership of specially-constructed target lists such as the Earth
Transit Zone. Such selections are entirely reasonable when the goal is to
maximise the a priori probability of detecting a genuine technosignature per
unit observing time, but they carry an unavoidable cost: they make it difficult
to convert a survey's results into a statistically clean statement about the
prevalence of technosignatures in the general stellar population, precisely
because the sample is not representative of that population
\citep{Margot2023}. A disc-hosting or exoplanet-hosting bias is particularly
worth flagging in the ALMA context: ALMA-detectable circumstellar and
protoplanetary discs are, almost without exception, orders of magnitude brighter
than the present-day Solar System's own zodiacal dust cloud would appear at
interstellar distances, so that samples selected \emph{because} they have an
ALMA-detected disc are biased towards young, dynamically active systems that are
not necessarily more (and could plausibly be less) likely to host a mature,
Earth-like technological civilisation than an old, quiescent, disc-poor star of
the same spectral type.

The present survey is explicitly designed to avoid this selection effect. We
construct our sample by identifying every star within 20\,pc of the Sun that
happens to fall within the field of view of at least one public ALMA
observation -- irrespective of whether that observation was proposed to study a
disc, a stellar flare, an unrelated background source, or anything else -- and
irrespective of the target star's own spectral type, age, multiplicity, or known
planet-hosting status. The star need not be at the pointing centre of the
observation, provided it lies within the usable primary-beam field of view (see
\S\ref{sec:sample}). This is, in effect, a volume-limited ``bycatch'' census of
the ALMA archive: rather than sampling stars selected for astrophysical interest
in the calibrator fields as in \citet{Mason2024}, we sample every nearby star
that ALMA happens to have observed for any purpose whatsoever, within a fixed
physical volume. We recognise, and quantify explicitly in Paper~II, that this
approach is still subject to a residual, second-order selection effect -- ALMA
archival coverage of the immediate solar neighbourhood is itself not spatially
uniform, since some regions and spectral types are more likely to have been
proposed for ALMA time than others -- but this bias operates on the presence of
\emph{any} ALMA coverage of a star, not on the astrophysical properties that are
the actual subject of a technosignature search, and can therefore be
characterised and reported as a well-defined completeness function of the
20\,pc volume, rather than an implicit and unquantified bias in the target
selection itself.

\subsection{Prospects for next-generation facilities}
\label{sec:future}

The coming decade will see a substantial expansion of technosignature search
capability. The Square Kilometre Array \citep[SKA;][]{SKA2015,SKA2026review},
now under construction in Australia (SKA-Low) and South Africa (SKA-Mid), will
combine wide instantaneous bandwidth, high angular resolution, and an
architecture well suited to commensal, multi-beam technosignature searching,
with projected sensitivity sufficient to detect Earth-analogue-level radar and
communication leakage from significantly greater distances than is possible
today. The next-generation Very Large Array \citep[ngVLA;][]{ngVLASETI2022} is
specifically being designed with technosignature science as one of its
commensal use cases, with proposed antenna configurations optimised to allow of
order 64 simultaneous high-spectral-resolution commensal beams. Existing
commensal systems at MeerKAT, the VLA (COSMIC), and the Murchison Widefield
Array are already providing a preview of the observing modes these larger
facilities will employ at scale. None of these facilities, however, offer ALMA's
combination of angular resolution, spectral resolution and raw collecting area
above $\sim$\,50\,GHz in the near term, so a systematic mm/submm technosignature
census -- of the kind begun here -- is likely to remain a uniquely ALMA-enabled
measurement for some years yet, and represents a natural complement to the
lower-frequency, larger-sample surveys that will be enabled by SKA and ngVLA.

\section{Sample selection}
\label{sec:sample}

Our sample comprises every star within 20\,pc of the Sun (based on Gaia~DR3
parallaxes) with at least one public ALMA observation, of any band, science
category, or proposal type, in which the star's true (proper-motion-corrected)
position at the epoch of observation falls within the field of view of at least
one execution block -- not necessarily at the observation's pointing centre,
since a substantial fraction of ALMA mosaics and wide single-pointing
observations serendipitously cover more than one catalogued star. We deliberately
do not weight target selection by the presence of a circumstellar disc, known
exoplanets, spectral type, or any other astrophysical property beyond distance
and archival data availability; disc presence and planet-hosting status are
retained purely as descriptive metadata for the final sample, not as
selection or ranking criteria. At the time of writing this identifies 120 unique
target stars, spanning spectral types from early A (e.g.\ Sirius~B, $\alpha$~CMa~B)
through to the latest, faintest M~dwarfs (e.g.\ Proxima Centauri, TRAPPIST-1,
Barnard's Star, Teegarden's Star), ranked by increasing distance from the Sun so
that the nearest, most sensitively-constrained systems are processed first. The
underlying pipeline is built to extend this volume-limited approach to 30, 40 or
50\,pc as a natural next step, at the cost of a rapidly increasing target list
and a correspondingly reduced fraction with existing ALMA coverage, since the
archive's areal completeness at fixed exposure declines with distance.

\section{Data and methodology}
\label{sec:method}

For each target we mine the public ALMA science archive for all calibrated (or,
where necessary, locally recalibrated from raw data) execution blocks covering
the star's position, apply an efficiency-bounded download and processing budget
per target to remain within the resource constraints of the shared compute
facility used for this work, and search the resulting visibility data in two
complementary ways: (i) a narrowband, Doppler-drift-corrected search for
technosignature candidates, following the general methodology of
\citet{Mason2024} and \citet{Enriquez2017}, converting any non-detection into an
equivalent isotropic radiated power (EIRP) limit via
$\mathrm{EIRP_{min}} = 4\pi d^{2} S_{\rm min} \Delta\nu$; and (ii) a
line-excluded continuum search for anomalous broad-band emission at the star's
position, with an accompanying epoch-to-epoch variability check wherever
multiple execution blocks exist. Stellar positions are propagated to the epoch
of each observation using Gaia~DR3 proper motions and parallaxes. The
Doppler-drift search range applied to each target is set to the larger of (a) a
generic frequency-scaled default following \citet{Mason2024}'s adoption of a
literature fiducial with a threefold safety margin, and (b) where the target
hosts one or more confirmed exoplanets, an explicit calculation of the maximum
line-of-sight orbital acceleration (and hence maximum plausible transmitter
drift rate) of the innermost known planet, so that a transmitter co-located with
a known, short-period planet cannot be missed simply because its orbital motion
drifts it out of an insufficiently wide search window. Full details of the
pipeline, its validation, and the specific corrections applied for
interferometric mosaic fields and for spectral-line contamination of continuum
measurements, will be given in Paper~II together with the resulting sample-wide
detection limits.

\section{Scope of Paper II}
\label{sec:futurework}

Paper~II of this series will present: the completed narrowband and continuum
technosignature search of the full 20\,pc sample described in
\S\ref{sec:sample}; per-target and aggregate EIRP and continuum flux-density
upper limits, tabulated together with the exact frequency ranges searched in
each case; any candidate signals identified, together with their vetting against
known interference and astrophysical contaminants; a discussion of the sample's
completeness as a function of spectral type and distance within 20\,pc; and an
assessment, following the statistical frameworks of \citet{Wright2018},
\citet{Sheikh2020} and \citet{Margot2023}, of what fraction of the local stellar
population our non-detections allow us to exclude as hosting a detectable
millimetre/submillimetre transmitter of Earth-2026-equivalent power. We will
also discuss prospects for extending the survey to 30, 40 and 50\,pc, and for
applying the anomaly-detection and machine-learning techniques reviewed in
\S\ref{sec:ml} to the same underlying dataset as a complementary search
methodology to the classical narrowband/continuum approach adopted here.

\section*{Acknowledgements}
This paper makes use of the following ALMA data: [to be completed in Paper~II
with the full list of ALMA project codes used]. ALMA is a partnership of ESO
(representing its member states), NSF (USA) and NINS (Japan), together with NRC
(Canada), MOST and ASIAA (Taiwan), and KASI (Republic of Korea), in cooperation
with the Republic of Chile. The Joint ALMA Observatory is operated by ESO,
AUI/NRAO and NAOJ. This work has made use of data from the European Space
Agency (ESA) mission {\it Gaia}, processed by the {\it Gaia} Data Processing and
Analysis Consortium (DPAC).

\section*{Data Availability}
The reduced data products underlying this paper, together with the full
target-selection and analysis pipeline, are being made available at
\url{https://github.com/Tilanthi/SETI} as the survey progresses. Raw ALMA
visibility data are publicly available from the ALMA Science Archive at
\url{https://almascience.org}.

\begin{thebibliography}{99}

\bibitem[Backus \& Project Phoenix Team(2002)]{Backus2002} Backus P.~R., Project Phoenix Team, 2002, in Bulletin of the American Astronomical Society. p.\ 1173

\bibitem[Choza et al.(2024)]{Choza2024} Choza C., et al., 2024, AJ, 168, 254

\bibitem[Cocconi \& Morrison(1959)]{CocconiMorrison1959} Cocconi G., Morrison P., 1959, Nature, 184, 844

\bibitem[Czech et al.(2021)]{Czech2021} Czech D., et al., 2021, PASP, 133, 064502

\bibitem[Czech et al.(2026)]{BLMeerKAT2026} Czech D.~J., et al., 2026, arXiv:2607.23651

\bibitem[Drake(1960)]{Drake1960} Drake F.~D., 1960, Physics Today, 13, 40

\bibitem[Drake(1974)]{Drake1974} Drake F.~D., 1974, in Sagan C., ed., Communication with Extraterrestrial Intelligence (CETI). MIT Press

\bibitem[Drake(1986)]{Drake1986} Drake F.~D., 1986, in Applications of Digital Image Processing IX. SPIE

\bibitem[Enriquez \& Price(2019)]{Enriquez2019} Enriquez E., Price D., 2019, turboSETI, Astrophysics Source Code Library, record ascl:1906.006

\bibitem[Enriquez et al.(2017)]{Enriquez2017} Enriquez J.~E., et al., 2017, ApJ, 849, 104

\bibitem[Gajjar et al.(2019)]{Gajjar2019EarthTransitZone} Gajjar V., et al., 2019, in Bulletin of the American Astronomical Society. AAS Meeting \#233

\bibitem[Garrett et al.(2026)]{BRaTs2026} Garrett M.~A., et al., 2026, MNRAS, accepted (arXiv:2605.10212)

\bibitem[Isaacson et al.(2017)]{BLoverview} Isaacson H., et al., 2017, PASP, 129, 054501

\bibitem[Jacobson-Bell et al.(2025)]{GLOBULAR2025} Jacobson-Bell K., et al., 2025, AJ, 169, 233

\bibitem[Johnson et al.(2023)]{LOFARdualsite2023} Johnson O.~A., et al., 2023, AJ, 166, 193

\bibitem[Lacki et al.(2021)]{GBTarchiveML2025} Lacki B.~C., et al., 2021, arXiv:2101.02244

\bibitem[Li et al.(2026)]{FAST3IATLAS2025} Li J.-K., Tao Z.-Z., Zhang T.-J., 2026, arXiv:2603.19023

\bibitem[Ma et al.(2023)]{Ma2023} Ma P.~X., et al., 2023, Nature Astronomy, 7, 492

\bibitem[Margot et al.(2023)]{Margot2023} Margot J.-L., et al., 2023, AJ, 166, 189

\bibitem[Mason et al.(2024)]{Mason2024} Mason L.~A., Garrett M.~A., Wandia K., Siemion A.~P.~V., 2024, MNRAS, 536, 2127

\bibitem[Morrison et al.(2023)]{Morrison2023} Morrison I.~S., et al., 2023, PASA, 40, e038

\bibitem[Ng et al.(2022)]{ngVLASETI2022} Ng C., et al., 2022, AJ, 164, 205

\bibitem[Poznanski et al.(2025)]{AnomalyParkesGBT2025} Poznanski D., et al., 2025, AJ, in press (arXiv:2505.03927)

\bibitem[Price et al.(2020)]{Price2020} Price D.~C., et al., 2020, AJ, 159, 86

\bibitem[Sheikh(2020)]{Sheikh2020} Sheikh S.~Z., 2020, International Journal of Astrobiology, 19, 237

\bibitem[Sheikh et al.(2025a)]{EarthDetectingEarth2025} Sheikh S.~Z., et al., 2025, AJ, 169, 108

\bibitem[Sheikh et al.(2025b)]{ATAISO2025} Sheikh S.~Z., et al., 2025, arXiv:2512.18142 (``A Search for Radio Technosignatures from Interstellar Object 3I/ATLAS with the Allen Telescope Array'')

\bibitem[Siemion et al.(2015)]{SKA2015} Siemion A.~P.~V., et al., 2015, arXiv:1412.4867

\bibitem[Tarter(2001)]{Tarter2001} Tarter J., 2001, ARA\&A, 39, 511

\bibitem[Tremblay et al.(2024)]{Tremblay2024} Tremblay C.~D., et al., 2024, PASA, 41, e004

\bibitem[Tremblay et al.(2026)]{SKA2026review} Tremblay C.~D., et al., 2026, in Advancing Astrophysics with the SKA II (AASKAII), arXiv:2606.27565

\bibitem[Tusay et al.(2024)]{ATATRAPPIST2024} Tusay N., Sheikh S.~Z., Sneed E.~L., Farah W., Pollak A.~W., Cruz L.~F., Siemion A., DeBoer D.~R., Wright J.~T., 2024, AJ, 168, 283

\bibitem[Tusay et al.(2025)]{TESSeclipsing2025} Tusay N., et al., 2025, arXiv:2506.13459

\bibitem[Valente et al.(2025)]{SRT2025} Valente G., et al., 2025, Acta Astronautica, in press

\bibitem[Vidal et al.(2026)]{SETIreview2025} Vidal C., et al., 2026, arXiv:2605.21093 (``The Search for Technosignatures: a Review of Possibilities'')

\bibitem[Von Korff et al.(2025)]{SETIhome2025} Von Korff J., et al., 2025, AJ, 170, 112

\bibitem[Wright et al.(2018)]{Wright2018} Wright J.~T., Kanodia S., Lubar E.~G., 2018, AJ, 156, 260

\bibitem[Zhang et al.(2020)]{Zhang2020} Zhang Z.-S., et al., 2020, ApJ, 891, 174

\end{thebibliography}

\end{document}
